% texto de orientacao em cinza; \mostrarorientacoesfalse tira todos do pdf
\newif\ifmostrarorientacoes
\mostrarorientacoestrue
\newcommand{\orientacao}[1]{\ifmostrarorientacoes{\color{gray}\itshape #1\par}\fi}

% dado ainda nao definido, em vermelho para nao passar despercebido
\newcommand{\pendente}[1]{\textcolor{red}{#1}}

% tela de imagens/prototipos, com moldura porque o fundo se confunde com a pagina
% uso: \prototipo{arquivo}{legenda}, referenciado com \ref{fig:prot-arquivo}
\newcommand{\prototipo}[2]{%
  \begin{figure}[H]
    \centering
    \caption{#2}
    \label{fig:prot-#1}
    \setlength{\fboxsep}{0pt}%
    \fbox{\includegraphics[width=\dimexpr\linewidth-2\fboxrule\relax]{imagens/prototipos/#1.jpg}}
    \fonte{Elaborado pelos autores.}
  \end{figure}%
}

% espaco reservado para diagrama ainda nao desenhado
% uso: \diagrama{legenda}{rotulo}, referenciado com \ref{fig:rotulo}
% quando a imagem existir em imagens/, trocar pelo ambiente figure com \includegraphics
\newcommand{\diagrama}[2]{%
  \begin{figure}[htbp]
    \centering
    \caption{#1}
    \label{fig:#2}
    \fbox{\parbox[c][5cm][c]{0.8\textwidth}{\centering\color{gray}\itshape #1 (a elaborar)}}
    \fonte{Elaborado pelos autores.}
  \end{figure}%
}
